% ----------------------------------------------------------
% 9. Projeto de Dados
% ----------------------------------------------------------
\chapter{Projeto de Dados}
\label{cap:projeto-dados}

% ----------------------------------------------------------
\section{Modelo de dados do projeto}
% ----------------------------------------------------------

\obs{Colocar aqui o diagrama que representa o modelo de dados do projeto: o DER do seu
projeto (DER: Diagrama Entidade Relacionamento), ou o esquema de dados NOSQL etc. Ele
pode ser elaborado em ferramentas de modelagem UML ou gerados pelo ambiente de
desenvolvimento. (Deve existir em todos os projetos que envolvem banco de dados). O
exemplo abaixo é de um DER para vocês lembrarem como é.}

\begin{figure}[htb]
	\centering
	\caption[Exemplo de DER]{Exemplo de DER}
	\label{fig:exemplo-der}
	\includegraphics[width=\textwidth]{exemplo-der}
\end{figure}

% ----------------------------------------------------------
\section{Dicionário de Dados Completo}
% ----------------------------------------------------------

\obs{Faça um dicionário completo. O exemplo abaixo é para vocês lembrarem como é.}

% ---
% O documento original pede que as páginas do dicionário de dados fiquem em PAISAGEM.
% Isso é feito aqui pelo ambiente landscape (pacote pdflscape).
% ---
\begin{landscape}

\footnotesize

\noindent
\begin{tabular}{|p{4.0cm}|p{12.0cm}|p{5.0cm}|}
\hline
\textbf{ENTIDADE} & & \\
\hline
\textbf{Nome:} & \multicolumn{2}{p{17.4cm}|}{Tabela de Curso} \\
\hline
\textbf{Nome Abreviado:} & \multicolumn{2}{p{17.4cm}|}{TBCURSO} \\
\hline
\textbf{Índices:} & \textbf{Nome do Índice} & \textbf{Atributos} \\
\hline
 & IND\_TBCURSO\_NMCURSO & NMCURSO \\
\hline
\textbf{Descrição:} & \multicolumn{2}{p{17.4cm}|}{Tabela que contém todos os cursos
oferecidos pela Universidade Miranda} \\
\hline
\end{tabular}

\vspace{0.4cm}

\noindent\textbf{Atributos:}

\noindent
\begin{tabular}{|p{3.6cm}|c|c|p{7.0cm}|p{2.8cm}|p{3.0cm}|}
\hline
\textbf{Nome do Atributo} & \textbf{Tipo} & \textbf{Tamanho} & \textbf{Descrição} &
\textbf{Máscara} & \textbf{Regra Validação} \\
\hline
\textbf{\underline{CDCURSO}} & N & 5,0 & Código do Curso & 99.999 & >0 \\
\hline
NMCURSO & T & 50 & Nome do Curso & & Obrigatório \\
\hline
CDAREA & N & 5,0 & Código da Área & 99.999 & >0 \\
\hline
DTINCREG & D & 10 & Data Inclusão Registro & 99/99/9999 & < = date ( ) \\
\hline
CDUSUINC & N & 5,0 & Código Usuário Incluiu Registro & 99.999 & >0 \\
\hline
DTALTREC & D & 8 & Data Alteração do Registro & 99/99/9999 & < = date ( ) \\
\hline
CDUSUALT & N & 5,0 & Código Usuário Alterou Registro & 99.999 & >0 \\
\hline
\end{tabular}

\vspace{1.0cm}

\noindent
\begin{tabular}{|p{4.0cm}|p{12.0cm}|p{5.0cm}|}
\hline
\textbf{ENTIDADE} & & \\
\hline
\textbf{Nome:} & \multicolumn{2}{p{17.4cm}|}{Cadastro de Aluno} \\
\hline
\textbf{Nome Abreviado:} & \multicolumn{2}{p{17.4cm}|}{TBALUNO} \\
\hline
\textbf{Índices:} & \textbf{Nome do Índice} & \textbf{Atributos} \\
\hline
 & IND\_TBALUNO\_NMALUNO & NMALUNO \\
\hline
 & IND\_TBALUNO\_CPF & CPF \\
\hline
\textbf{Descrição:} & \multicolumn{2}{p{17.4cm}|}{Tabela que contém todos os alunos
matriculados na Universidade Miranda} \\
\hline
\end{tabular}

\vspace{0.4cm}

\noindent\textbf{Atributos:}

\noindent
\begin{tabular}{|p{3.6cm}|c|c|p{7.0cm}|p{2.8cm}|p{3.0cm}|}
\hline
\textbf{Nome do Atributo} & \textbf{Tipo} & \textbf{Tamanho} & \textbf{Descrição} &
\textbf{Máscara} & \textbf{Regra Validação} \\
\hline
\textbf{\underline{CDALUNO}} & N & 5,0 & Código do Aluno & 99.999 & >0 \\
\hline
NMALUNO & T & 50 & Nome do Aluno & & Obrigatório \\
\hline
SEXO & T & 1 & Sexo & & \\
\hline
CDCURSO & N & 5,0 & Código do Curso & 99.999 & >0 \\
\hline
CDCIDADENASC & N & 5,0 & Código Cidade Nascimento & 99.999 & >0 \\
\hline
CDCIDADERES & N & 5,0 & Código Cidade Residência & 99.999 & >0 \\
\hline
EDFONE & T & 10 & Telefone & (99)9999-9999 & \\
\hline
CPF & N & 11 & CPF do Aluno & 999.999.999-99 & obrigatório \\
\hline
DTINCREG & D & 8 & Data Inclusão Registro & 99/99/9999 & < = date ( ) \\
\hline
CDUSUINC & N & 5,0 & Código Usuário Incluiu Registro & 99.999 & >0 \\
\hline
DTALTREC & D & 8 & Data Alteração do Registro & 99/99/9999 & < = date ( ) \\
\hline
CDUSUALT & N & 5,0 & Código Usuário Alterou Registro & 99.999 & >0 \\
\hline
\end{tabular}

\normalsize
\end{landscape}
